% ============================================================================
% ESQUEMA (banca): o laço do aprendizado ativo com acervo fixo (pool-based),
% na formulação de Settles e no arcabouço da tese (sêxtupla
% A = <L, U, theta, S, O, B>), com a seleção EM LOTE destacada.
% Encomenda do autor (tarefa 0036 do principal) para o §2.2
% (2-fundam, sec:formalismo) — FUNDAMENTO: oráculo genérico, sem
% antecipar a maquinaria do Cap. 3.
% Uso previsto (o principal gateia a inserção; a banca NÃO edita o Cap. 2):
%   após o Algoritmo alg:active_learning, dentro de figure:
%   % ============================================================================
% ESQUEMA (banca): o laço do aprendizado ativo com acervo fixo (pool-based),
% na formulação de Settles e no arcabouço da tese (sêxtupla
% A = <L, U, theta, S, O, B>), com a seleção EM LOTE destacada.
% Encomenda do autor (tarefa 0036 do principal) para o §2.2
% (2-fundam, sec:formalismo) — FUNDAMENTO: oráculo genérico, sem
% antecipar a maquinaria do Cap. 3.
% Uso previsto (o principal gateia a inserção; a banca NÃO edita o Cap. 2):
%   após o Algoritmo alg:active_learning, dentro de figure:
%   % ============================================================================
% ESQUEMA (banca): o laço do aprendizado ativo com acervo fixo (pool-based),
% na formulação de Settles e no arcabouço da tese (sêxtupla
% A = <L, U, theta, S, O, B>), com a seleção EM LOTE destacada.
% Encomenda do autor (tarefa 0036 do principal) para o §2.2
% (2-fundam, sec:formalismo) — FUNDAMENTO: oráculo genérico, sem
% antecipar a maquinaria do Cap. 3.
% Uso previsto (o principal gateia a inserção; a banca NÃO edita o Cap. 2):
%   após o Algoritmo alg:active_learning, dentro de figure:
%   % ============================================================================
% ESQUEMA (banca): o laço do aprendizado ativo com acervo fixo (pool-based),
% na formulação de Settles e no arcabouço da tese (sêxtupla
% A = <L, U, theta, S, O, B>), com a seleção EM LOTE destacada.
% Encomenda do autor (tarefa 0036 do principal) para o §2.2
% (2-fundam, sec:formalismo) — FUNDAMENTO: oráculo genérico, sem
% antecipar a maquinaria do Cap. 3.
% Uso previsto (o principal gateia a inserção; a banca NÃO edita o Cap. 2):
%   após o Algoritmo alg:active_learning, dentro de figure:
%   \input{3-metodo/esquemas-propostos/esq-laco-aprendizado-ativo}
% Notação idêntica à Tabela tab:active-learning-components e ao Algoritmo 1.
% SEM números medidos, SEM códigos internos, SEM caminhos, SEM travessões.
% Uso standalone (pré-visualização):
%   \documentclass[tikz,border=6pt,12pt]{standalone}
%   \usetikzlibrary{arrows.meta, positioning}
%   \begin{document}\input{esq-laco-aprendizado-ativo.tex}\end{document}
% GEOMETRIA CALIBRADA PARA CORPO 12 (ppginf/book 12pt, textwidth 16cm).
% Requer apenas arrows.meta e positioning (já em packages.tex). Legível em
% P&B: estado/dados = cinza claro; oráculo e saída = cinza médio;
% operações = branco; saída do laço = tracejado.
% ============================================================================
\begin{tikzpicture}[
    x=1cm, y=1cm,
    caixa/.style={draw, rounded corners=2pt, align=center, text width=42mm,
                  minimum height=13mm, inner sep=3pt, font=\footnotesize},
    dado/.style={caixa, fill=gray!8},
    oraculo/.style={caixa, fill=gray!14},
    saida/.style={caixa, fill=gray!14},
    seta/.style={-{Stealth[length=2.6mm]}, thick},
    rot/.style={font=\scriptsize, align=center, inner sep=2pt}
]
    % ------------------ o acervo, fora do laço -----------------------------
    \node[dado] (U) at (10.8,2.8)
        {acervo não rotulado $U$\\(\textit{pool} fixo, disponível\\desde o início)};

    % ------------------ o laço (anel de cinco nós) -------------------------
    \node[dado] (L) at (0,0)
        {conjunto rotulado $L_t$\\(inicia como $L_0$)};
    \node[caixa] (treino) at (5.4,0)
        {treinamento do classificador:\\$\theta_t \leftarrow \mathrm{Learn}(L_t)$};
    \node[caixa] (selecao) at (10.8,0)
        {seleção de um \textbf{lote}:\\$Q_t \leftarrow S(U_t,\theta_t,L_t)$\\
         (modo em lote, \textit{batch-mode})};
    \node[oraculo] (orac) at (10.8,-3.4)
        {\textbf{oráculo} $O$ rotula\\cada instância do lote:\\$y \leftarrow O(x)$};
    \node[caixa, text width=46mm] (atualiza) at (5.4,-3.4)
        {atualização:\\
         $L_{t+1} \leftarrow L_t \cup \{(x,y)\}$\\
         $U_{t+1} \leftarrow U_t \setminus Q_t$;\;
         \mbox{$B_{t+1} \leftarrow B_t - |Q_t|$}};

    % ------------------ saída do laço (abaixo do anel) ---------------------
    \node[saida, text width=46mm] (fim) at (5.4,-6.3)
        {\textbf{parada}: orçamento $B$\\esgotado ou acervo exausto;\\
         devolve $\theta_{\mathrm{final}}$ treinado em $L_T$};

    % ------------------ setas ----------------------------------------------
    \draw[seta] (U) -- (selecao);
    \draw[seta] (L) -- (treino);
    \draw[seta] (treino) -- (selecao);
    \draw[seta] (selecao) -- (orac);
    \draw[seta] (orac) -- (atualiza);
    % o retorno fecha o anel por caminho reto: a atualização devolve os
    % pares rotulados a L (a saída de parada fica abaixo, fora do corredor)
    \draw[seta] (atualiza.west) -| (L.south);
    \node[rot, anchor=south] at (1.55,-3.3)
        {repete enquanto\\houver orçamento};
    \draw[seta, dashed] (atualiza) -- (fim);
\end{tikzpicture}

% Notação idêntica à Tabela tab:active-learning-components e ao Algoritmo 1.
% SEM números medidos, SEM códigos internos, SEM caminhos, SEM travessões.
% Uso standalone (pré-visualização):
%   \documentclass[tikz,border=6pt,12pt]{standalone}
%   \usetikzlibrary{arrows.meta, positioning}
%   \begin{document}% ============================================================================
% ESQUEMA (banca): o laço do aprendizado ativo com acervo fixo (pool-based),
% na formulação de Settles e no arcabouço da tese (sêxtupla
% A = <L, U, theta, S, O, B>), com a seleção EM LOTE destacada.
% Encomenda do autor (tarefa 0036 do principal) para o §2.2
% (2-fundam, sec:formalismo) — FUNDAMENTO: oráculo genérico, sem
% antecipar a maquinaria do Cap. 3.
% Uso previsto (o principal gateia a inserção; a banca NÃO edita o Cap. 2):
%   após o Algoritmo alg:active_learning, dentro de figure:
%   \input{3-metodo/esquemas-propostos/esq-laco-aprendizado-ativo}
% Notação idêntica à Tabela tab:active-learning-components e ao Algoritmo 1.
% SEM números medidos, SEM códigos internos, SEM caminhos, SEM travessões.
% Uso standalone (pré-visualização):
%   \documentclass[tikz,border=6pt,12pt]{standalone}
%   \usetikzlibrary{arrows.meta, positioning}
%   \begin{document}\input{esq-laco-aprendizado-ativo.tex}\end{document}
% GEOMETRIA CALIBRADA PARA CORPO 12 (ppginf/book 12pt, textwidth 16cm).
% Requer apenas arrows.meta e positioning (já em packages.tex). Legível em
% P&B: estado/dados = cinza claro; oráculo e saída = cinza médio;
% operações = branco; saída do laço = tracejado.
% ============================================================================
\begin{tikzpicture}[
    x=1cm, y=1cm,
    caixa/.style={draw, rounded corners=2pt, align=center, text width=42mm,
                  minimum height=13mm, inner sep=3pt, font=\footnotesize},
    dado/.style={caixa, fill=gray!8},
    oraculo/.style={caixa, fill=gray!14},
    saida/.style={caixa, fill=gray!14},
    seta/.style={-{Stealth[length=2.6mm]}, thick},
    rot/.style={font=\scriptsize, align=center, inner sep=2pt}
]
    % ------------------ o acervo, fora do laço -----------------------------
    \node[dado] (U) at (10.8,2.8)
        {acervo não rotulado $U$\\(\textit{pool} fixo, disponível\\desde o início)};

    % ------------------ o laço (anel de cinco nós) -------------------------
    \node[dado] (L) at (0,0)
        {conjunto rotulado $L_t$\\(inicia como $L_0$)};
    \node[caixa] (treino) at (5.4,0)
        {treinamento do classificador:\\$\theta_t \leftarrow \mathrm{Learn}(L_t)$};
    \node[caixa] (selecao) at (10.8,0)
        {seleção de um \textbf{lote}:\\$Q_t \leftarrow S(U_t,\theta_t,L_t)$\\
         (modo em lote, \textit{batch-mode})};
    \node[oraculo] (orac) at (10.8,-3.4)
        {\textbf{oráculo} $O$ rotula\\cada instância do lote:\\$y \leftarrow O(x)$};
    \node[caixa, text width=46mm] (atualiza) at (5.4,-3.4)
        {atualização:\\
         $L_{t+1} \leftarrow L_t \cup \{(x,y)\}$\\
         $U_{t+1} \leftarrow U_t \setminus Q_t$;\;
         \mbox{$B_{t+1} \leftarrow B_t - |Q_t|$}};

    % ------------------ saída do laço (abaixo do anel) ---------------------
    \node[saida, text width=46mm] (fim) at (5.4,-6.3)
        {\textbf{parada}: orçamento $B$\\esgotado ou acervo exausto;\\
         devolve $\theta_{\mathrm{final}}$ treinado em $L_T$};

    % ------------------ setas ----------------------------------------------
    \draw[seta] (U) -- (selecao);
    \draw[seta] (L) -- (treino);
    \draw[seta] (treino) -- (selecao);
    \draw[seta] (selecao) -- (orac);
    \draw[seta] (orac) -- (atualiza);
    % o retorno fecha o anel por caminho reto: a atualização devolve os
    % pares rotulados a L (a saída de parada fica abaixo, fora do corredor)
    \draw[seta] (atualiza.west) -| (L.south);
    \node[rot, anchor=south] at (1.55,-3.3)
        {repete enquanto\\houver orçamento};
    \draw[seta, dashed] (atualiza) -- (fim);
\end{tikzpicture}
\end{document}
% GEOMETRIA CALIBRADA PARA CORPO 12 (ppginf/book 12pt, textwidth 16cm).
% Requer apenas arrows.meta e positioning (já em packages.tex). Legível em
% P&B: estado/dados = cinza claro; oráculo e saída = cinza médio;
% operações = branco; saída do laço = tracejado.
% ============================================================================
\begin{tikzpicture}[
    x=1cm, y=1cm,
    caixa/.style={draw, rounded corners=2pt, align=center, text width=42mm,
                  minimum height=13mm, inner sep=3pt, font=\footnotesize},
    dado/.style={caixa, fill=gray!8},
    oraculo/.style={caixa, fill=gray!14},
    saida/.style={caixa, fill=gray!14},
    seta/.style={-{Stealth[length=2.6mm]}, thick},
    rot/.style={font=\scriptsize, align=center, inner sep=2pt}
]
    % ------------------ o acervo, fora do laço -----------------------------
    \node[dado] (U) at (10.8,2.8)
        {acervo não rotulado $U$\\(\textit{pool} fixo, disponível\\desde o início)};

    % ------------------ o laço (anel de cinco nós) -------------------------
    \node[dado] (L) at (0,0)
        {conjunto rotulado $L_t$\\(inicia como $L_0$)};
    \node[caixa] (treino) at (5.4,0)
        {treinamento do classificador:\\$\theta_t \leftarrow \mathrm{Learn}(L_t)$};
    \node[caixa] (selecao) at (10.8,0)
        {seleção de um \textbf{lote}:\\$Q_t \leftarrow S(U_t,\theta_t,L_t)$\\
         (modo em lote, \textit{batch-mode})};
    \node[oraculo] (orac) at (10.8,-3.4)
        {\textbf{oráculo} $O$ rotula\\cada instância do lote:\\$y \leftarrow O(x)$};
    \node[caixa, text width=46mm] (atualiza) at (5.4,-3.4)
        {atualização:\\
         $L_{t+1} \leftarrow L_t \cup \{(x,y)\}$\\
         $U_{t+1} \leftarrow U_t \setminus Q_t$;\;
         \mbox{$B_{t+1} \leftarrow B_t - |Q_t|$}};

    % ------------------ saída do laço (abaixo do anel) ---------------------
    \node[saida, text width=46mm] (fim) at (5.4,-6.3)
        {\textbf{parada}: orçamento $B$\\esgotado ou acervo exausto;\\
         devolve $\theta_{\mathrm{final}}$ treinado em $L_T$};

    % ------------------ setas ----------------------------------------------
    \draw[seta] (U) -- (selecao);
    \draw[seta] (L) -- (treino);
    \draw[seta] (treino) -- (selecao);
    \draw[seta] (selecao) -- (orac);
    \draw[seta] (orac) -- (atualiza);
    % o retorno fecha o anel por caminho reto: a atualização devolve os
    % pares rotulados a L (a saída de parada fica abaixo, fora do corredor)
    \draw[seta] (atualiza.west) -| (L.south);
    \node[rot, anchor=south] at (1.55,-3.3)
        {repete enquanto\\houver orçamento};
    \draw[seta, dashed] (atualiza) -- (fim);
\end{tikzpicture}

% Notação idêntica à Tabela tab:active-learning-components e ao Algoritmo 1.
% SEM números medidos, SEM códigos internos, SEM caminhos, SEM travessões.
% Uso standalone (pré-visualização):
%   \documentclass[tikz,border=6pt,12pt]{standalone}
%   \usetikzlibrary{arrows.meta, positioning}
%   \begin{document}% ============================================================================
% ESQUEMA (banca): o laço do aprendizado ativo com acervo fixo (pool-based),
% na formulação de Settles e no arcabouço da tese (sêxtupla
% A = <L, U, theta, S, O, B>), com a seleção EM LOTE destacada.
% Encomenda do autor (tarefa 0036 do principal) para o §2.2
% (2-fundam, sec:formalismo) — FUNDAMENTO: oráculo genérico, sem
% antecipar a maquinaria do Cap. 3.
% Uso previsto (o principal gateia a inserção; a banca NÃO edita o Cap. 2):
%   após o Algoritmo alg:active_learning, dentro de figure:
%   % ============================================================================
% ESQUEMA (banca): o laço do aprendizado ativo com acervo fixo (pool-based),
% na formulação de Settles e no arcabouço da tese (sêxtupla
% A = <L, U, theta, S, O, B>), com a seleção EM LOTE destacada.
% Encomenda do autor (tarefa 0036 do principal) para o §2.2
% (2-fundam, sec:formalismo) — FUNDAMENTO: oráculo genérico, sem
% antecipar a maquinaria do Cap. 3.
% Uso previsto (o principal gateia a inserção; a banca NÃO edita o Cap. 2):
%   após o Algoritmo alg:active_learning, dentro de figure:
%   \input{3-metodo/esquemas-propostos/esq-laco-aprendizado-ativo}
% Notação idêntica à Tabela tab:active-learning-components e ao Algoritmo 1.
% SEM números medidos, SEM códigos internos, SEM caminhos, SEM travessões.
% Uso standalone (pré-visualização):
%   \documentclass[tikz,border=6pt,12pt]{standalone}
%   \usetikzlibrary{arrows.meta, positioning}
%   \begin{document}\input{esq-laco-aprendizado-ativo.tex}\end{document}
% GEOMETRIA CALIBRADA PARA CORPO 12 (ppginf/book 12pt, textwidth 16cm).
% Requer apenas arrows.meta e positioning (já em packages.tex). Legível em
% P&B: estado/dados = cinza claro; oráculo e saída = cinza médio;
% operações = branco; saída do laço = tracejado.
% ============================================================================
\begin{tikzpicture}[
    x=1cm, y=1cm,
    caixa/.style={draw, rounded corners=2pt, align=center, text width=42mm,
                  minimum height=13mm, inner sep=3pt, font=\footnotesize},
    dado/.style={caixa, fill=gray!8},
    oraculo/.style={caixa, fill=gray!14},
    saida/.style={caixa, fill=gray!14},
    seta/.style={-{Stealth[length=2.6mm]}, thick},
    rot/.style={font=\scriptsize, align=center, inner sep=2pt}
]
    % ------------------ o acervo, fora do laço -----------------------------
    \node[dado] (U) at (10.8,2.8)
        {acervo não rotulado $U$\\(\textit{pool} fixo, disponível\\desde o início)};

    % ------------------ o laço (anel de cinco nós) -------------------------
    \node[dado] (L) at (0,0)
        {conjunto rotulado $L_t$\\(inicia como $L_0$)};
    \node[caixa] (treino) at (5.4,0)
        {treinamento do classificador:\\$\theta_t \leftarrow \mathrm{Learn}(L_t)$};
    \node[caixa] (selecao) at (10.8,0)
        {seleção de um \textbf{lote}:\\$Q_t \leftarrow S(U_t,\theta_t,L_t)$\\
         (modo em lote, \textit{batch-mode})};
    \node[oraculo] (orac) at (10.8,-3.4)
        {\textbf{oráculo} $O$ rotula\\cada instância do lote:\\$y \leftarrow O(x)$};
    \node[caixa, text width=46mm] (atualiza) at (5.4,-3.4)
        {atualização:\\
         $L_{t+1} \leftarrow L_t \cup \{(x,y)\}$\\
         $U_{t+1} \leftarrow U_t \setminus Q_t$;\;
         \mbox{$B_{t+1} \leftarrow B_t - |Q_t|$}};

    % ------------------ saída do laço (abaixo do anel) ---------------------
    \node[saida, text width=46mm] (fim) at (5.4,-6.3)
        {\textbf{parada}: orçamento $B$\\esgotado ou acervo exausto;\\
         devolve $\theta_{\mathrm{final}}$ treinado em $L_T$};

    % ------------------ setas ----------------------------------------------
    \draw[seta] (U) -- (selecao);
    \draw[seta] (L) -- (treino);
    \draw[seta] (treino) -- (selecao);
    \draw[seta] (selecao) -- (orac);
    \draw[seta] (orac) -- (atualiza);
    % o retorno fecha o anel por caminho reto: a atualização devolve os
    % pares rotulados a L (a saída de parada fica abaixo, fora do corredor)
    \draw[seta] (atualiza.west) -| (L.south);
    \node[rot, anchor=south] at (1.55,-3.3)
        {repete enquanto\\houver orçamento};
    \draw[seta, dashed] (atualiza) -- (fim);
\end{tikzpicture}

% Notação idêntica à Tabela tab:active-learning-components e ao Algoritmo 1.
% SEM números medidos, SEM códigos internos, SEM caminhos, SEM travessões.
% Uso standalone (pré-visualização):
%   \documentclass[tikz,border=6pt,12pt]{standalone}
%   \usetikzlibrary{arrows.meta, positioning}
%   \begin{document}% ============================================================================
% ESQUEMA (banca): o laço do aprendizado ativo com acervo fixo (pool-based),
% na formulação de Settles e no arcabouço da tese (sêxtupla
% A = <L, U, theta, S, O, B>), com a seleção EM LOTE destacada.
% Encomenda do autor (tarefa 0036 do principal) para o §2.2
% (2-fundam, sec:formalismo) — FUNDAMENTO: oráculo genérico, sem
% antecipar a maquinaria do Cap. 3.
% Uso previsto (o principal gateia a inserção; a banca NÃO edita o Cap. 2):
%   após o Algoritmo alg:active_learning, dentro de figure:
%   \input{3-metodo/esquemas-propostos/esq-laco-aprendizado-ativo}
% Notação idêntica à Tabela tab:active-learning-components e ao Algoritmo 1.
% SEM números medidos, SEM códigos internos, SEM caminhos, SEM travessões.
% Uso standalone (pré-visualização):
%   \documentclass[tikz,border=6pt,12pt]{standalone}
%   \usetikzlibrary{arrows.meta, positioning}
%   \begin{document}\input{esq-laco-aprendizado-ativo.tex}\end{document}
% GEOMETRIA CALIBRADA PARA CORPO 12 (ppginf/book 12pt, textwidth 16cm).
% Requer apenas arrows.meta e positioning (já em packages.tex). Legível em
% P&B: estado/dados = cinza claro; oráculo e saída = cinza médio;
% operações = branco; saída do laço = tracejado.
% ============================================================================
\begin{tikzpicture}[
    x=1cm, y=1cm,
    caixa/.style={draw, rounded corners=2pt, align=center, text width=42mm,
                  minimum height=13mm, inner sep=3pt, font=\footnotesize},
    dado/.style={caixa, fill=gray!8},
    oraculo/.style={caixa, fill=gray!14},
    saida/.style={caixa, fill=gray!14},
    seta/.style={-{Stealth[length=2.6mm]}, thick},
    rot/.style={font=\scriptsize, align=center, inner sep=2pt}
]
    % ------------------ o acervo, fora do laço -----------------------------
    \node[dado] (U) at (10.8,2.8)
        {acervo não rotulado $U$\\(\textit{pool} fixo, disponível\\desde o início)};

    % ------------------ o laço (anel de cinco nós) -------------------------
    \node[dado] (L) at (0,0)
        {conjunto rotulado $L_t$\\(inicia como $L_0$)};
    \node[caixa] (treino) at (5.4,0)
        {treinamento do classificador:\\$\theta_t \leftarrow \mathrm{Learn}(L_t)$};
    \node[caixa] (selecao) at (10.8,0)
        {seleção de um \textbf{lote}:\\$Q_t \leftarrow S(U_t,\theta_t,L_t)$\\
         (modo em lote, \textit{batch-mode})};
    \node[oraculo] (orac) at (10.8,-3.4)
        {\textbf{oráculo} $O$ rotula\\cada instância do lote:\\$y \leftarrow O(x)$};
    \node[caixa, text width=46mm] (atualiza) at (5.4,-3.4)
        {atualização:\\
         $L_{t+1} \leftarrow L_t \cup \{(x,y)\}$\\
         $U_{t+1} \leftarrow U_t \setminus Q_t$;\;
         \mbox{$B_{t+1} \leftarrow B_t - |Q_t|$}};

    % ------------------ saída do laço (abaixo do anel) ---------------------
    \node[saida, text width=46mm] (fim) at (5.4,-6.3)
        {\textbf{parada}: orçamento $B$\\esgotado ou acervo exausto;\\
         devolve $\theta_{\mathrm{final}}$ treinado em $L_T$};

    % ------------------ setas ----------------------------------------------
    \draw[seta] (U) -- (selecao);
    \draw[seta] (L) -- (treino);
    \draw[seta] (treino) -- (selecao);
    \draw[seta] (selecao) -- (orac);
    \draw[seta] (orac) -- (atualiza);
    % o retorno fecha o anel por caminho reto: a atualização devolve os
    % pares rotulados a L (a saída de parada fica abaixo, fora do corredor)
    \draw[seta] (atualiza.west) -| (L.south);
    \node[rot, anchor=south] at (1.55,-3.3)
        {repete enquanto\\houver orçamento};
    \draw[seta, dashed] (atualiza) -- (fim);
\end{tikzpicture}
\end{document}
% GEOMETRIA CALIBRADA PARA CORPO 12 (ppginf/book 12pt, textwidth 16cm).
% Requer apenas arrows.meta e positioning (já em packages.tex). Legível em
% P&B: estado/dados = cinza claro; oráculo e saída = cinza médio;
% operações = branco; saída do laço = tracejado.
% ============================================================================
\begin{tikzpicture}[
    x=1cm, y=1cm,
    caixa/.style={draw, rounded corners=2pt, align=center, text width=42mm,
                  minimum height=13mm, inner sep=3pt, font=\footnotesize},
    dado/.style={caixa, fill=gray!8},
    oraculo/.style={caixa, fill=gray!14},
    saida/.style={caixa, fill=gray!14},
    seta/.style={-{Stealth[length=2.6mm]}, thick},
    rot/.style={font=\scriptsize, align=center, inner sep=2pt}
]
    % ------------------ o acervo, fora do laço -----------------------------
    \node[dado] (U) at (10.8,2.8)
        {acervo não rotulado $U$\\(\textit{pool} fixo, disponível\\desde o início)};

    % ------------------ o laço (anel de cinco nós) -------------------------
    \node[dado] (L) at (0,0)
        {conjunto rotulado $L_t$\\(inicia como $L_0$)};
    \node[caixa] (treino) at (5.4,0)
        {treinamento do classificador:\\$\theta_t \leftarrow \mathrm{Learn}(L_t)$};
    \node[caixa] (selecao) at (10.8,0)
        {seleção de um \textbf{lote}:\\$Q_t \leftarrow S(U_t,\theta_t,L_t)$\\
         (modo em lote, \textit{batch-mode})};
    \node[oraculo] (orac) at (10.8,-3.4)
        {\textbf{oráculo} $O$ rotula\\cada instância do lote:\\$y \leftarrow O(x)$};
    \node[caixa, text width=46mm] (atualiza) at (5.4,-3.4)
        {atualização:\\
         $L_{t+1} \leftarrow L_t \cup \{(x,y)\}$\\
         $U_{t+1} \leftarrow U_t \setminus Q_t$;\;
         \mbox{$B_{t+1} \leftarrow B_t - |Q_t|$}};

    % ------------------ saída do laço (abaixo do anel) ---------------------
    \node[saida, text width=46mm] (fim) at (5.4,-6.3)
        {\textbf{parada}: orçamento $B$\\esgotado ou acervo exausto;\\
         devolve $\theta_{\mathrm{final}}$ treinado em $L_T$};

    % ------------------ setas ----------------------------------------------
    \draw[seta] (U) -- (selecao);
    \draw[seta] (L) -- (treino);
    \draw[seta] (treino) -- (selecao);
    \draw[seta] (selecao) -- (orac);
    \draw[seta] (orac) -- (atualiza);
    % o retorno fecha o anel por caminho reto: a atualização devolve os
    % pares rotulados a L (a saída de parada fica abaixo, fora do corredor)
    \draw[seta] (atualiza.west) -| (L.south);
    \node[rot, anchor=south] at (1.55,-3.3)
        {repete enquanto\\houver orçamento};
    \draw[seta, dashed] (atualiza) -- (fim);
\end{tikzpicture}
\end{document}
% GEOMETRIA CALIBRADA PARA CORPO 12 (ppginf/book 12pt, textwidth 16cm).
% Requer apenas arrows.meta e positioning (já em packages.tex). Legível em
% P&B: estado/dados = cinza claro; oráculo e saída = cinza médio;
% operações = branco; saída do laço = tracejado.
% ============================================================================
\begin{tikzpicture}[
    x=1cm, y=1cm,
    caixa/.style={draw, rounded corners=2pt, align=center, text width=42mm,
                  minimum height=13mm, inner sep=3pt, font=\footnotesize},
    dado/.style={caixa, fill=gray!8},
    oraculo/.style={caixa, fill=gray!14},
    saida/.style={caixa, fill=gray!14},
    seta/.style={-{Stealth[length=2.6mm]}, thick},
    rot/.style={font=\scriptsize, align=center, inner sep=2pt}
]
    % ------------------ o acervo, fora do laço -----------------------------
    \node[dado] (U) at (10.8,2.8)
        {acervo não rotulado $U$\\(\textit{pool} fixo, disponível\\desde o início)};

    % ------------------ o laço (anel de cinco nós) -------------------------
    \node[dado] (L) at (0,0)
        {conjunto rotulado $L_t$\\(inicia como $L_0$)};
    \node[caixa] (treino) at (5.4,0)
        {treinamento do classificador:\\$\theta_t \leftarrow \mathrm{Learn}(L_t)$};
    \node[caixa] (selecao) at (10.8,0)
        {seleção de um \textbf{lote}:\\$Q_t \leftarrow S(U_t,\theta_t,L_t)$\\
         (modo em lote, \textit{batch-mode})};
    \node[oraculo] (orac) at (10.8,-3.4)
        {\textbf{oráculo} $O$ rotula\\cada instância do lote:\\$y \leftarrow O(x)$};
    \node[caixa, text width=46mm] (atualiza) at (5.4,-3.4)
        {atualização:\\
         $L_{t+1} \leftarrow L_t \cup \{(x,y)\}$\\
         $U_{t+1} \leftarrow U_t \setminus Q_t$;\;
         \mbox{$B_{t+1} \leftarrow B_t - |Q_t|$}};

    % ------------------ saída do laço (abaixo do anel) ---------------------
    \node[saida, text width=46mm] (fim) at (5.4,-6.3)
        {\textbf{parada}: orçamento $B$\\esgotado ou acervo exausto;\\
         devolve $\theta_{\mathrm{final}}$ treinado em $L_T$};

    % ------------------ setas ----------------------------------------------
    \draw[seta] (U) -- (selecao);
    \draw[seta] (L) -- (treino);
    \draw[seta] (treino) -- (selecao);
    \draw[seta] (selecao) -- (orac);
    \draw[seta] (orac) -- (atualiza);
    % o retorno fecha o anel por caminho reto: a atualização devolve os
    % pares rotulados a L (a saída de parada fica abaixo, fora do corredor)
    \draw[seta] (atualiza.west) -| (L.south);
    \node[rot, anchor=south] at (1.55,-3.3)
        {repete enquanto\\houver orçamento};
    \draw[seta, dashed] (atualiza) -- (fim);
\end{tikzpicture}

% Notação idêntica à Tabela tab:active-learning-components e ao Algoritmo 1.
% SEM números medidos, SEM códigos internos, SEM caminhos, SEM travessões.
% Uso standalone (pré-visualização):
%   \documentclass[tikz,border=6pt,12pt]{standalone}
%   \usetikzlibrary{arrows.meta, positioning}
%   \begin{document}% ============================================================================
% ESQUEMA (banca): o laço do aprendizado ativo com acervo fixo (pool-based),
% na formulação de Settles e no arcabouço da tese (sêxtupla
% A = <L, U, theta, S, O, B>), com a seleção EM LOTE destacada.
% Encomenda do autor (tarefa 0036 do principal) para o §2.2
% (2-fundam, sec:formalismo) — FUNDAMENTO: oráculo genérico, sem
% antecipar a maquinaria do Cap. 3.
% Uso previsto (o principal gateia a inserção; a banca NÃO edita o Cap. 2):
%   após o Algoritmo alg:active_learning, dentro de figure:
%   % ============================================================================
% ESQUEMA (banca): o laço do aprendizado ativo com acervo fixo (pool-based),
% na formulação de Settles e no arcabouço da tese (sêxtupla
% A = <L, U, theta, S, O, B>), com a seleção EM LOTE destacada.
% Encomenda do autor (tarefa 0036 do principal) para o §2.2
% (2-fundam, sec:formalismo) — FUNDAMENTO: oráculo genérico, sem
% antecipar a maquinaria do Cap. 3.
% Uso previsto (o principal gateia a inserção; a banca NÃO edita o Cap. 2):
%   após o Algoritmo alg:active_learning, dentro de figure:
%   % ============================================================================
% ESQUEMA (banca): o laço do aprendizado ativo com acervo fixo (pool-based),
% na formulação de Settles e no arcabouço da tese (sêxtupla
% A = <L, U, theta, S, O, B>), com a seleção EM LOTE destacada.
% Encomenda do autor (tarefa 0036 do principal) para o §2.2
% (2-fundam, sec:formalismo) — FUNDAMENTO: oráculo genérico, sem
% antecipar a maquinaria do Cap. 3.
% Uso previsto (o principal gateia a inserção; a banca NÃO edita o Cap. 2):
%   após o Algoritmo alg:active_learning, dentro de figure:
%   \input{3-metodo/esquemas-propostos/esq-laco-aprendizado-ativo}
% Notação idêntica à Tabela tab:active-learning-components e ao Algoritmo 1.
% SEM números medidos, SEM códigos internos, SEM caminhos, SEM travessões.
% Uso standalone (pré-visualização):
%   \documentclass[tikz,border=6pt,12pt]{standalone}
%   \usetikzlibrary{arrows.meta, positioning}
%   \begin{document}\input{esq-laco-aprendizado-ativo.tex}\end{document}
% GEOMETRIA CALIBRADA PARA CORPO 12 (ppginf/book 12pt, textwidth 16cm).
% Requer apenas arrows.meta e positioning (já em packages.tex). Legível em
% P&B: estado/dados = cinza claro; oráculo e saída = cinza médio;
% operações = branco; saída do laço = tracejado.
% ============================================================================
\begin{tikzpicture}[
    x=1cm, y=1cm,
    caixa/.style={draw, rounded corners=2pt, align=center, text width=42mm,
                  minimum height=13mm, inner sep=3pt, font=\footnotesize},
    dado/.style={caixa, fill=gray!8},
    oraculo/.style={caixa, fill=gray!14},
    saida/.style={caixa, fill=gray!14},
    seta/.style={-{Stealth[length=2.6mm]}, thick},
    rot/.style={font=\scriptsize, align=center, inner sep=2pt}
]
    % ------------------ o acervo, fora do laço -----------------------------
    \node[dado] (U) at (10.8,2.8)
        {acervo não rotulado $U$\\(\textit{pool} fixo, disponível\\desde o início)};

    % ------------------ o laço (anel de cinco nós) -------------------------
    \node[dado] (L) at (0,0)
        {conjunto rotulado $L_t$\\(inicia como $L_0$)};
    \node[caixa] (treino) at (5.4,0)
        {treinamento do classificador:\\$\theta_t \leftarrow \mathrm{Learn}(L_t)$};
    \node[caixa] (selecao) at (10.8,0)
        {seleção de um \textbf{lote}:\\$Q_t \leftarrow S(U_t,\theta_t,L_t)$\\
         (modo em lote, \textit{batch-mode})};
    \node[oraculo] (orac) at (10.8,-3.4)
        {\textbf{oráculo} $O$ rotula\\cada instância do lote:\\$y \leftarrow O(x)$};
    \node[caixa, text width=46mm] (atualiza) at (5.4,-3.4)
        {atualização:\\
         $L_{t+1} \leftarrow L_t \cup \{(x,y)\}$\\
         $U_{t+1} \leftarrow U_t \setminus Q_t$;\;
         \mbox{$B_{t+1} \leftarrow B_t - |Q_t|$}};

    % ------------------ saída do laço (abaixo do anel) ---------------------
    \node[saida, text width=46mm] (fim) at (5.4,-6.3)
        {\textbf{parada}: orçamento $B$\\esgotado ou acervo exausto;\\
         devolve $\theta_{\mathrm{final}}$ treinado em $L_T$};

    % ------------------ setas ----------------------------------------------
    \draw[seta] (U) -- (selecao);
    \draw[seta] (L) -- (treino);
    \draw[seta] (treino) -- (selecao);
    \draw[seta] (selecao) -- (orac);
    \draw[seta] (orac) -- (atualiza);
    % o retorno fecha o anel por caminho reto: a atualização devolve os
    % pares rotulados a L (a saída de parada fica abaixo, fora do corredor)
    \draw[seta] (atualiza.west) -| (L.south);
    \node[rot, anchor=south] at (1.55,-3.3)
        {repete enquanto\\houver orçamento};
    \draw[seta, dashed] (atualiza) -- (fim);
\end{tikzpicture}

% Notação idêntica à Tabela tab:active-learning-components e ao Algoritmo 1.
% SEM números medidos, SEM códigos internos, SEM caminhos, SEM travessões.
% Uso standalone (pré-visualização):
%   \documentclass[tikz,border=6pt,12pt]{standalone}
%   \usetikzlibrary{arrows.meta, positioning}
%   \begin{document}% ============================================================================
% ESQUEMA (banca): o laço do aprendizado ativo com acervo fixo (pool-based),
% na formulação de Settles e no arcabouço da tese (sêxtupla
% A = <L, U, theta, S, O, B>), com a seleção EM LOTE destacada.
% Encomenda do autor (tarefa 0036 do principal) para o §2.2
% (2-fundam, sec:formalismo) — FUNDAMENTO: oráculo genérico, sem
% antecipar a maquinaria do Cap. 3.
% Uso previsto (o principal gateia a inserção; a banca NÃO edita o Cap. 2):
%   após o Algoritmo alg:active_learning, dentro de figure:
%   \input{3-metodo/esquemas-propostos/esq-laco-aprendizado-ativo}
% Notação idêntica à Tabela tab:active-learning-components e ao Algoritmo 1.
% SEM números medidos, SEM códigos internos, SEM caminhos, SEM travessões.
% Uso standalone (pré-visualização):
%   \documentclass[tikz,border=6pt,12pt]{standalone}
%   \usetikzlibrary{arrows.meta, positioning}
%   \begin{document}\input{esq-laco-aprendizado-ativo.tex}\end{document}
% GEOMETRIA CALIBRADA PARA CORPO 12 (ppginf/book 12pt, textwidth 16cm).
% Requer apenas arrows.meta e positioning (já em packages.tex). Legível em
% P&B: estado/dados = cinza claro; oráculo e saída = cinza médio;
% operações = branco; saída do laço = tracejado.
% ============================================================================
\begin{tikzpicture}[
    x=1cm, y=1cm,
    caixa/.style={draw, rounded corners=2pt, align=center, text width=42mm,
                  minimum height=13mm, inner sep=3pt, font=\footnotesize},
    dado/.style={caixa, fill=gray!8},
    oraculo/.style={caixa, fill=gray!14},
    saida/.style={caixa, fill=gray!14},
    seta/.style={-{Stealth[length=2.6mm]}, thick},
    rot/.style={font=\scriptsize, align=center, inner sep=2pt}
]
    % ------------------ o acervo, fora do laço -----------------------------
    \node[dado] (U) at (10.8,2.8)
        {acervo não rotulado $U$\\(\textit{pool} fixo, disponível\\desde o início)};

    % ------------------ o laço (anel de cinco nós) -------------------------
    \node[dado] (L) at (0,0)
        {conjunto rotulado $L_t$\\(inicia como $L_0$)};
    \node[caixa] (treino) at (5.4,0)
        {treinamento do classificador:\\$\theta_t \leftarrow \mathrm{Learn}(L_t)$};
    \node[caixa] (selecao) at (10.8,0)
        {seleção de um \textbf{lote}:\\$Q_t \leftarrow S(U_t,\theta_t,L_t)$\\
         (modo em lote, \textit{batch-mode})};
    \node[oraculo] (orac) at (10.8,-3.4)
        {\textbf{oráculo} $O$ rotula\\cada instância do lote:\\$y \leftarrow O(x)$};
    \node[caixa, text width=46mm] (atualiza) at (5.4,-3.4)
        {atualização:\\
         $L_{t+1} \leftarrow L_t \cup \{(x,y)\}$\\
         $U_{t+1} \leftarrow U_t \setminus Q_t$;\;
         \mbox{$B_{t+1} \leftarrow B_t - |Q_t|$}};

    % ------------------ saída do laço (abaixo do anel) ---------------------
    \node[saida, text width=46mm] (fim) at (5.4,-6.3)
        {\textbf{parada}: orçamento $B$\\esgotado ou acervo exausto;\\
         devolve $\theta_{\mathrm{final}}$ treinado em $L_T$};

    % ------------------ setas ----------------------------------------------
    \draw[seta] (U) -- (selecao);
    \draw[seta] (L) -- (treino);
    \draw[seta] (treino) -- (selecao);
    \draw[seta] (selecao) -- (orac);
    \draw[seta] (orac) -- (atualiza);
    % o retorno fecha o anel por caminho reto: a atualização devolve os
    % pares rotulados a L (a saída de parada fica abaixo, fora do corredor)
    \draw[seta] (atualiza.west) -| (L.south);
    \node[rot, anchor=south] at (1.55,-3.3)
        {repete enquanto\\houver orçamento};
    \draw[seta, dashed] (atualiza) -- (fim);
\end{tikzpicture}
\end{document}
% GEOMETRIA CALIBRADA PARA CORPO 12 (ppginf/book 12pt, textwidth 16cm).
% Requer apenas arrows.meta e positioning (já em packages.tex). Legível em
% P&B: estado/dados = cinza claro; oráculo e saída = cinza médio;
% operações = branco; saída do laço = tracejado.
% ============================================================================
\begin{tikzpicture}[
    x=1cm, y=1cm,
    caixa/.style={draw, rounded corners=2pt, align=center, text width=42mm,
                  minimum height=13mm, inner sep=3pt, font=\footnotesize},
    dado/.style={caixa, fill=gray!8},
    oraculo/.style={caixa, fill=gray!14},
    saida/.style={caixa, fill=gray!14},
    seta/.style={-{Stealth[length=2.6mm]}, thick},
    rot/.style={font=\scriptsize, align=center, inner sep=2pt}
]
    % ------------------ o acervo, fora do laço -----------------------------
    \node[dado] (U) at (10.8,2.8)
        {acervo não rotulado $U$\\(\textit{pool} fixo, disponível\\desde o início)};

    % ------------------ o laço (anel de cinco nós) -------------------------
    \node[dado] (L) at (0,0)
        {conjunto rotulado $L_t$\\(inicia como $L_0$)};
    \node[caixa] (treino) at (5.4,0)
        {treinamento do classificador:\\$\theta_t \leftarrow \mathrm{Learn}(L_t)$};
    \node[caixa] (selecao) at (10.8,0)
        {seleção de um \textbf{lote}:\\$Q_t \leftarrow S(U_t,\theta_t,L_t)$\\
         (modo em lote, \textit{batch-mode})};
    \node[oraculo] (orac) at (10.8,-3.4)
        {\textbf{oráculo} $O$ rotula\\cada instância do lote:\\$y \leftarrow O(x)$};
    \node[caixa, text width=46mm] (atualiza) at (5.4,-3.4)
        {atualização:\\
         $L_{t+1} \leftarrow L_t \cup \{(x,y)\}$\\
         $U_{t+1} \leftarrow U_t \setminus Q_t$;\;
         \mbox{$B_{t+1} \leftarrow B_t - |Q_t|$}};

    % ------------------ saída do laço (abaixo do anel) ---------------------
    \node[saida, text width=46mm] (fim) at (5.4,-6.3)
        {\textbf{parada}: orçamento $B$\\esgotado ou acervo exausto;\\
         devolve $\theta_{\mathrm{final}}$ treinado em $L_T$};

    % ------------------ setas ----------------------------------------------
    \draw[seta] (U) -- (selecao);
    \draw[seta] (L) -- (treino);
    \draw[seta] (treino) -- (selecao);
    \draw[seta] (selecao) -- (orac);
    \draw[seta] (orac) -- (atualiza);
    % o retorno fecha o anel por caminho reto: a atualização devolve os
    % pares rotulados a L (a saída de parada fica abaixo, fora do corredor)
    \draw[seta] (atualiza.west) -| (L.south);
    \node[rot, anchor=south] at (1.55,-3.3)
        {repete enquanto\\houver orçamento};
    \draw[seta, dashed] (atualiza) -- (fim);
\end{tikzpicture}

% Notação idêntica à Tabela tab:active-learning-components e ao Algoritmo 1.
% SEM números medidos, SEM códigos internos, SEM caminhos, SEM travessões.
% Uso standalone (pré-visualização):
%   \documentclass[tikz,border=6pt,12pt]{standalone}
%   \usetikzlibrary{arrows.meta, positioning}
%   \begin{document}% ============================================================================
% ESQUEMA (banca): o laço do aprendizado ativo com acervo fixo (pool-based),
% na formulação de Settles e no arcabouço da tese (sêxtupla
% A = <L, U, theta, S, O, B>), com a seleção EM LOTE destacada.
% Encomenda do autor (tarefa 0036 do principal) para o §2.2
% (2-fundam, sec:formalismo) — FUNDAMENTO: oráculo genérico, sem
% antecipar a maquinaria do Cap. 3.
% Uso previsto (o principal gateia a inserção; a banca NÃO edita o Cap. 2):
%   após o Algoritmo alg:active_learning, dentro de figure:
%   % ============================================================================
% ESQUEMA (banca): o laço do aprendizado ativo com acervo fixo (pool-based),
% na formulação de Settles e no arcabouço da tese (sêxtupla
% A = <L, U, theta, S, O, B>), com a seleção EM LOTE destacada.
% Encomenda do autor (tarefa 0036 do principal) para o §2.2
% (2-fundam, sec:formalismo) — FUNDAMENTO: oráculo genérico, sem
% antecipar a maquinaria do Cap. 3.
% Uso previsto (o principal gateia a inserção; a banca NÃO edita o Cap. 2):
%   após o Algoritmo alg:active_learning, dentro de figure:
%   \input{3-metodo/esquemas-propostos/esq-laco-aprendizado-ativo}
% Notação idêntica à Tabela tab:active-learning-components e ao Algoritmo 1.
% SEM números medidos, SEM códigos internos, SEM caminhos, SEM travessões.
% Uso standalone (pré-visualização):
%   \documentclass[tikz,border=6pt,12pt]{standalone}
%   \usetikzlibrary{arrows.meta, positioning}
%   \begin{document}\input{esq-laco-aprendizado-ativo.tex}\end{document}
% GEOMETRIA CALIBRADA PARA CORPO 12 (ppginf/book 12pt, textwidth 16cm).
% Requer apenas arrows.meta e positioning (já em packages.tex). Legível em
% P&B: estado/dados = cinza claro; oráculo e saída = cinza médio;
% operações = branco; saída do laço = tracejado.
% ============================================================================
\begin{tikzpicture}[
    x=1cm, y=1cm,
    caixa/.style={draw, rounded corners=2pt, align=center, text width=42mm,
                  minimum height=13mm, inner sep=3pt, font=\footnotesize},
    dado/.style={caixa, fill=gray!8},
    oraculo/.style={caixa, fill=gray!14},
    saida/.style={caixa, fill=gray!14},
    seta/.style={-{Stealth[length=2.6mm]}, thick},
    rot/.style={font=\scriptsize, align=center, inner sep=2pt}
]
    % ------------------ o acervo, fora do laço -----------------------------
    \node[dado] (U) at (10.8,2.8)
        {acervo não rotulado $U$\\(\textit{pool} fixo, disponível\\desde o início)};

    % ------------------ o laço (anel de cinco nós) -------------------------
    \node[dado] (L) at (0,0)
        {conjunto rotulado $L_t$\\(inicia como $L_0$)};
    \node[caixa] (treino) at (5.4,0)
        {treinamento do classificador:\\$\theta_t \leftarrow \mathrm{Learn}(L_t)$};
    \node[caixa] (selecao) at (10.8,0)
        {seleção de um \textbf{lote}:\\$Q_t \leftarrow S(U_t,\theta_t,L_t)$\\
         (modo em lote, \textit{batch-mode})};
    \node[oraculo] (orac) at (10.8,-3.4)
        {\textbf{oráculo} $O$ rotula\\cada instância do lote:\\$y \leftarrow O(x)$};
    \node[caixa, text width=46mm] (atualiza) at (5.4,-3.4)
        {atualização:\\
         $L_{t+1} \leftarrow L_t \cup \{(x,y)\}$\\
         $U_{t+1} \leftarrow U_t \setminus Q_t$;\;
         \mbox{$B_{t+1} \leftarrow B_t - |Q_t|$}};

    % ------------------ saída do laço (abaixo do anel) ---------------------
    \node[saida, text width=46mm] (fim) at (5.4,-6.3)
        {\textbf{parada}: orçamento $B$\\esgotado ou acervo exausto;\\
         devolve $\theta_{\mathrm{final}}$ treinado em $L_T$};

    % ------------------ setas ----------------------------------------------
    \draw[seta] (U) -- (selecao);
    \draw[seta] (L) -- (treino);
    \draw[seta] (treino) -- (selecao);
    \draw[seta] (selecao) -- (orac);
    \draw[seta] (orac) -- (atualiza);
    % o retorno fecha o anel por caminho reto: a atualização devolve os
    % pares rotulados a L (a saída de parada fica abaixo, fora do corredor)
    \draw[seta] (atualiza.west) -| (L.south);
    \node[rot, anchor=south] at (1.55,-3.3)
        {repete enquanto\\houver orçamento};
    \draw[seta, dashed] (atualiza) -- (fim);
\end{tikzpicture}

% Notação idêntica à Tabela tab:active-learning-components e ao Algoritmo 1.
% SEM números medidos, SEM códigos internos, SEM caminhos, SEM travessões.
% Uso standalone (pré-visualização):
%   \documentclass[tikz,border=6pt,12pt]{standalone}
%   \usetikzlibrary{arrows.meta, positioning}
%   \begin{document}% ============================================================================
% ESQUEMA (banca): o laço do aprendizado ativo com acervo fixo (pool-based),
% na formulação de Settles e no arcabouço da tese (sêxtupla
% A = <L, U, theta, S, O, B>), com a seleção EM LOTE destacada.
% Encomenda do autor (tarefa 0036 do principal) para o §2.2
% (2-fundam, sec:formalismo) — FUNDAMENTO: oráculo genérico, sem
% antecipar a maquinaria do Cap. 3.
% Uso previsto (o principal gateia a inserção; a banca NÃO edita o Cap. 2):
%   após o Algoritmo alg:active_learning, dentro de figure:
%   \input{3-metodo/esquemas-propostos/esq-laco-aprendizado-ativo}
% Notação idêntica à Tabela tab:active-learning-components e ao Algoritmo 1.
% SEM números medidos, SEM códigos internos, SEM caminhos, SEM travessões.
% Uso standalone (pré-visualização):
%   \documentclass[tikz,border=6pt,12pt]{standalone}
%   \usetikzlibrary{arrows.meta, positioning}
%   \begin{document}\input{esq-laco-aprendizado-ativo.tex}\end{document}
% GEOMETRIA CALIBRADA PARA CORPO 12 (ppginf/book 12pt, textwidth 16cm).
% Requer apenas arrows.meta e positioning (já em packages.tex). Legível em
% P&B: estado/dados = cinza claro; oráculo e saída = cinza médio;
% operações = branco; saída do laço = tracejado.
% ============================================================================
\begin{tikzpicture}[
    x=1cm, y=1cm,
    caixa/.style={draw, rounded corners=2pt, align=center, text width=42mm,
                  minimum height=13mm, inner sep=3pt, font=\footnotesize},
    dado/.style={caixa, fill=gray!8},
    oraculo/.style={caixa, fill=gray!14},
    saida/.style={caixa, fill=gray!14},
    seta/.style={-{Stealth[length=2.6mm]}, thick},
    rot/.style={font=\scriptsize, align=center, inner sep=2pt}
]
    % ------------------ o acervo, fora do laço -----------------------------
    \node[dado] (U) at (10.8,2.8)
        {acervo não rotulado $U$\\(\textit{pool} fixo, disponível\\desde o início)};

    % ------------------ o laço (anel de cinco nós) -------------------------
    \node[dado] (L) at (0,0)
        {conjunto rotulado $L_t$\\(inicia como $L_0$)};
    \node[caixa] (treino) at (5.4,0)
        {treinamento do classificador:\\$\theta_t \leftarrow \mathrm{Learn}(L_t)$};
    \node[caixa] (selecao) at (10.8,0)
        {seleção de um \textbf{lote}:\\$Q_t \leftarrow S(U_t,\theta_t,L_t)$\\
         (modo em lote, \textit{batch-mode})};
    \node[oraculo] (orac) at (10.8,-3.4)
        {\textbf{oráculo} $O$ rotula\\cada instância do lote:\\$y \leftarrow O(x)$};
    \node[caixa, text width=46mm] (atualiza) at (5.4,-3.4)
        {atualização:\\
         $L_{t+1} \leftarrow L_t \cup \{(x,y)\}$\\
         $U_{t+1} \leftarrow U_t \setminus Q_t$;\;
         \mbox{$B_{t+1} \leftarrow B_t - |Q_t|$}};

    % ------------------ saída do laço (abaixo do anel) ---------------------
    \node[saida, text width=46mm] (fim) at (5.4,-6.3)
        {\textbf{parada}: orçamento $B$\\esgotado ou acervo exausto;\\
         devolve $\theta_{\mathrm{final}}$ treinado em $L_T$};

    % ------------------ setas ----------------------------------------------
    \draw[seta] (U) -- (selecao);
    \draw[seta] (L) -- (treino);
    \draw[seta] (treino) -- (selecao);
    \draw[seta] (selecao) -- (orac);
    \draw[seta] (orac) -- (atualiza);
    % o retorno fecha o anel por caminho reto: a atualização devolve os
    % pares rotulados a L (a saída de parada fica abaixo, fora do corredor)
    \draw[seta] (atualiza.west) -| (L.south);
    \node[rot, anchor=south] at (1.55,-3.3)
        {repete enquanto\\houver orçamento};
    \draw[seta, dashed] (atualiza) -- (fim);
\end{tikzpicture}
\end{document}
% GEOMETRIA CALIBRADA PARA CORPO 12 (ppginf/book 12pt, textwidth 16cm).
% Requer apenas arrows.meta e positioning (já em packages.tex). Legível em
% P&B: estado/dados = cinza claro; oráculo e saída = cinza médio;
% operações = branco; saída do laço = tracejado.
% ============================================================================
\begin{tikzpicture}[
    x=1cm, y=1cm,
    caixa/.style={draw, rounded corners=2pt, align=center, text width=42mm,
                  minimum height=13mm, inner sep=3pt, font=\footnotesize},
    dado/.style={caixa, fill=gray!8},
    oraculo/.style={caixa, fill=gray!14},
    saida/.style={caixa, fill=gray!14},
    seta/.style={-{Stealth[length=2.6mm]}, thick},
    rot/.style={font=\scriptsize, align=center, inner sep=2pt}
]
    % ------------------ o acervo, fora do laço -----------------------------
    \node[dado] (U) at (10.8,2.8)
        {acervo não rotulado $U$\\(\textit{pool} fixo, disponível\\desde o início)};

    % ------------------ o laço (anel de cinco nós) -------------------------
    \node[dado] (L) at (0,0)
        {conjunto rotulado $L_t$\\(inicia como $L_0$)};
    \node[caixa] (treino) at (5.4,0)
        {treinamento do classificador:\\$\theta_t \leftarrow \mathrm{Learn}(L_t)$};
    \node[caixa] (selecao) at (10.8,0)
        {seleção de um \textbf{lote}:\\$Q_t \leftarrow S(U_t,\theta_t,L_t)$\\
         (modo em lote, \textit{batch-mode})};
    \node[oraculo] (orac) at (10.8,-3.4)
        {\textbf{oráculo} $O$ rotula\\cada instância do lote:\\$y \leftarrow O(x)$};
    \node[caixa, text width=46mm] (atualiza) at (5.4,-3.4)
        {atualização:\\
         $L_{t+1} \leftarrow L_t \cup \{(x,y)\}$\\
         $U_{t+1} \leftarrow U_t \setminus Q_t$;\;
         \mbox{$B_{t+1} \leftarrow B_t - |Q_t|$}};

    % ------------------ saída do laço (abaixo do anel) ---------------------
    \node[saida, text width=46mm] (fim) at (5.4,-6.3)
        {\textbf{parada}: orçamento $B$\\esgotado ou acervo exausto;\\
         devolve $\theta_{\mathrm{final}}$ treinado em $L_T$};

    % ------------------ setas ----------------------------------------------
    \draw[seta] (U) -- (selecao);
    \draw[seta] (L) -- (treino);
    \draw[seta] (treino) -- (selecao);
    \draw[seta] (selecao) -- (orac);
    \draw[seta] (orac) -- (atualiza);
    % o retorno fecha o anel por caminho reto: a atualização devolve os
    % pares rotulados a L (a saída de parada fica abaixo, fora do corredor)
    \draw[seta] (atualiza.west) -| (L.south);
    \node[rot, anchor=south] at (1.55,-3.3)
        {repete enquanto\\houver orçamento};
    \draw[seta, dashed] (atualiza) -- (fim);
\end{tikzpicture}
\end{document}
% GEOMETRIA CALIBRADA PARA CORPO 12 (ppginf/book 12pt, textwidth 16cm).
% Requer apenas arrows.meta e positioning (já em packages.tex). Legível em
% P&B: estado/dados = cinza claro; oráculo e saída = cinza médio;
% operações = branco; saída do laço = tracejado.
% ============================================================================
\begin{tikzpicture}[
    x=1cm, y=1cm,
    caixa/.style={draw, rounded corners=2pt, align=center, text width=42mm,
                  minimum height=13mm, inner sep=3pt, font=\footnotesize},
    dado/.style={caixa, fill=gray!8},
    oraculo/.style={caixa, fill=gray!14},
    saida/.style={caixa, fill=gray!14},
    seta/.style={-{Stealth[length=2.6mm]}, thick},
    rot/.style={font=\scriptsize, align=center, inner sep=2pt}
]
    % ------------------ o acervo, fora do laço -----------------------------
    \node[dado] (U) at (10.8,2.8)
        {acervo não rotulado $U$\\(\textit{pool} fixo, disponível\\desde o início)};

    % ------------------ o laço (anel de cinco nós) -------------------------
    \node[dado] (L) at (0,0)
        {conjunto rotulado $L_t$\\(inicia como $L_0$)};
    \node[caixa] (treino) at (5.4,0)
        {treinamento do classificador:\\$\theta_t \leftarrow \mathrm{Learn}(L_t)$};
    \node[caixa] (selecao) at (10.8,0)
        {seleção de um \textbf{lote}:\\$Q_t \leftarrow S(U_t,\theta_t,L_t)$\\
         (modo em lote, \textit{batch-mode})};
    \node[oraculo] (orac) at (10.8,-3.4)
        {\textbf{oráculo} $O$ rotula\\cada instância do lote:\\$y \leftarrow O(x)$};
    \node[caixa, text width=46mm] (atualiza) at (5.4,-3.4)
        {atualização:\\
         $L_{t+1} \leftarrow L_t \cup \{(x,y)\}$\\
         $U_{t+1} \leftarrow U_t \setminus Q_t$;\;
         \mbox{$B_{t+1} \leftarrow B_t - |Q_t|$}};

    % ------------------ saída do laço (abaixo do anel) ---------------------
    \node[saida, text width=46mm] (fim) at (5.4,-6.3)
        {\textbf{parada}: orçamento $B$\\esgotado ou acervo exausto;\\
         devolve $\theta_{\mathrm{final}}$ treinado em $L_T$};

    % ------------------ setas ----------------------------------------------
    \draw[seta] (U) -- (selecao);
    \draw[seta] (L) -- (treino);
    \draw[seta] (treino) -- (selecao);
    \draw[seta] (selecao) -- (orac);
    \draw[seta] (orac) -- (atualiza);
    % o retorno fecha o anel por caminho reto: a atualização devolve os
    % pares rotulados a L (a saída de parada fica abaixo, fora do corredor)
    \draw[seta] (atualiza.west) -| (L.south);
    \node[rot, anchor=south] at (1.55,-3.3)
        {repete enquanto\\houver orçamento};
    \draw[seta, dashed] (atualiza) -- (fim);
\end{tikzpicture}
\end{document}
% GEOMETRIA CALIBRADA PARA CORPO 12 (ppginf/book 12pt, textwidth 16cm).
% Requer apenas arrows.meta e positioning (já em packages.tex). Legível em
% P&B: estado/dados = cinza claro; oráculo e saída = cinza médio;
% operações = branco; saída do laço = tracejado.
% ============================================================================
\begin{tikzpicture}[
    x=1cm, y=1cm,
    caixa/.style={draw, rounded corners=2pt, align=center, text width=42mm,
                  minimum height=13mm, inner sep=3pt, font=\footnotesize},
    dado/.style={caixa, fill=gray!8},
    oraculo/.style={caixa, fill=gray!14},
    saida/.style={caixa, fill=gray!14},
    seta/.style={-{Stealth[length=2.6mm]}, thick},
    rot/.style={font=\scriptsize, align=center, inner sep=2pt}
]
    % ------------------ o acervo, fora do laço -----------------------------
    \node[dado] (U) at (10.8,2.8)
        {acervo não rotulado $U$\\(\textit{pool} fixo, disponível\\desde o início)};

    % ------------------ o laço (anel de cinco nós) -------------------------
    \node[dado] (L) at (0,0)
        {conjunto rotulado $L_t$\\(inicia como $L_0$)};
    \node[caixa] (treino) at (5.4,0)
        {treinamento do classificador:\\$\theta_t \leftarrow \mathrm{Learn}(L_t)$};
    \node[caixa] (selecao) at (10.8,0)
        {seleção de um \textbf{lote}:\\$Q_t \leftarrow S(U_t,\theta_t,L_t)$\\
         (modo em lote, \textit{batch-mode})};
    \node[oraculo] (orac) at (10.8,-3.4)
        {\textbf{oráculo} $O$ rotula\\cada instância do lote:\\$y \leftarrow O(x)$};
    \node[caixa, text width=46mm] (atualiza) at (5.4,-3.4)
        {atualização:\\
         $L_{t+1} \leftarrow L_t \cup \{(x,y)\}$\\
         $U_{t+1} \leftarrow U_t \setminus Q_t$;\;
         \mbox{$B_{t+1} \leftarrow B_t - |Q_t|$}};

    % ------------------ saída do laço (abaixo do anel) ---------------------
    \node[saida, text width=46mm] (fim) at (5.4,-6.3)
        {\textbf{parada}: orçamento $B$\\esgotado ou acervo exausto;\\
         devolve $\theta_{\mathrm{final}}$ treinado em $L_T$};

    % ------------------ setas ----------------------------------------------
    \draw[seta] (U) -- (selecao);
    \draw[seta] (L) -- (treino);
    \draw[seta] (treino) -- (selecao);
    \draw[seta] (selecao) -- (orac);
    \draw[seta] (orac) -- (atualiza);
    % o retorno fecha o anel por caminho reto: a atualização devolve os
    % pares rotulados a L (a saída de parada fica abaixo, fora do corredor)
    \draw[seta] (atualiza.west) -| (L.south);
    \node[rot, anchor=south] at (1.55,-3.3)
        {repete enquanto\\houver orçamento};
    \draw[seta, dashed] (atualiza) -- (fim);
\end{tikzpicture}
